\documentclass[12pt]{article}
\PassOptionsToPackage{table}{xcolor}
\usepackage[utf8]{inputenc}
\usepackage[T1]{fontenc}
\usepackage[italian]{babel}
\usepackage{multicol}
\usepackage{enumitem}
\usepackage{array}
\usepackage{longtable}
\usepackage{booktabs}
\usepackage{xcolor}
\usepackage{graphicx}
\usepackage{tikz}
\usetikzlibrary{
    shapes.geometric,
    arrows.meta,
    positioning,
    calc,
    patterns,
    decorations.pathmorphing
}
\usepackage[margin=2cm]{geometry}
\usepackage{amsmath,amssymb}
\setlength{\parindent}{0pt}
\definecolor{headercolor}{RGB}{220,220,220}
\newcounter{quesito}
\setcounter{quesito}{0}
\newcommand{\nuovoquesito}{\stepcounter{quesito}\textbf{Domanda \arabic{quesito}.}}
\newcommand{\itemdomanda}[2]{%
    \par\noindent
    \begin{minipage}[t]{\linewidth}
        \nuovoquesito \\ #1
        \begin{enumerate}[label=(\Alph*), nosep, leftmargin=*]
            #2
        \end{enumerate}
    \end{minipage}
    \vspace{10pt}
}
\begin{document}
\section*{Simulazione}
\hfill Data: 17 apr 2026

\itemdomanda{Se dopodomani sarà giovedì, che giorno era ieri?}{
    \item Lunedì
    \item Martedì
    \item Domenica
    \item Mercoledì
    \item Sabato
}

\itemdomanda{Un mattone pesa un chilo più mezzo mattone. Quanto pesa un mattone?}{
    \item \SI{1.5}{kg}
    \item \SI{2}{kg}
    \item \SI{2.5}{kg}
    \item \SI{3}{kg}
    \item \SI{1}{kg}
}

\itemdomanda{Cinque amici sono in fila: Anna è davanti a Luca, Sara è dietro a Luca ma davanti a Paolo. Se Marco è il primo della fila, chi è l'ultimo?}{
    \item Anna
    \item Luca
    \item Sara
    \item Paolo
    \item Marco
}

\itemdomanda{Ci sono cinque persone con diverse situazioni patrimoniali. Oronzo è più ricco di Rocco, le cui ricchezze sono più modeste di quelle di Silvio, e quest'ultimo a sua volta è più danaroso di Piero. Quirino è meno benestante di Piero, ma più agiato di Oronzo. Chi è il terzo in ordine di ricchezza?}{
    \item Piero
    \item Rocco
    \item Oronzo
    \item Silvio
    \item Quirino
}

\itemdomanda{Quale numero va inserito al centro?
\begin{center}
\begin{tikzpicture}[scale=0.5]
    \draw (0,0) grid (3,1);
    \node at (0.5,0.5) {14};
    \node at (1.5,0.5) {?};
    \node at (2.5,0.5) {28};
    \node[below] at (1.5,0) {\small (Relazione: $x = (A+B)/2$)};
\end{tikzpicture}
\end{center}}{
    \item 21
    \item 20
    \item 19
    \item 22
    \item 42
}

\itemdomanda{Completare la serie: $2, 6, 12, 20, 30, \dots$}{
    \item 36
    \item 40
    \item 42
    \item 50
    \item 44
}

\itemdomanda{Se un liquido scorre in una condotta con restringimento, quale di questi grafici rappresenta la portata in massa?}{
    \item \begin{tikzpicture}[scale=0.5, baseline=(current bounding box.center)]
    % Assi con stile locale
    \draw[-{Stealth[scale=1.2]}, thick] (0,0) -- (3.5,0) node[below, pos=0.8] {posizione};
    \draw[-{Stealth[scale=1.2]}, thick] (0,0) -- (0,3.5) node[left, rotate=90, pos=0.8, anchor=south] {portata};
    
    % Curva Parabolica (Concavità verso l'alto)
    % Funzione: y = 0.5*(x-1.5)^2 + 0.8
    \draw[thick, samples=100, domain=0.2:2.8] plot (\x, {0.5*(\x-1.5)*(\x-1.5) + 0.8});
\end{tikzpicture}
    \item \begin{tikzpicture}[scale=0.5, baseline=(current bounding box.center)]
    % Assi
    \draw[-{Stealth[scale=1.2]}, thick] (0,0) -- (3.5,0) node[below, pos=0.8] {posizione};
    \draw[-{Stealth[scale=1.2]}, thick] (0,0) -- (0,3.5) node[left, rotate=90, pos=0.8, anchor=south] {portata};
    
    % Linea Orizzontale (Costante)
    \draw[thick] (0,1.8) -- (3,1.8);
\end{tikzpicture}
    \item \begin{tikzpicture}[scale=0.5, baseline=(current bounding box.center)]
    % Assi
    \draw[-{Stealth[scale=1.2]}, thick] (0,0) -- (3.5,0) node[below, pos=0.8] {posizione};
    \draw[-{Stealth[scale=1.2]}, thick] (0,0) -- (0,3.5) node[left, rotate=90, pos=0.8, anchor=south] {portata};
    
    % Curva Parabolica (Concavità verso il basso)
    % Funzione: y = -0.5*(x-1.5)^2 + 2.2
    \draw[thick, samples=100, domain=0.1:2.9] plot (\x, {-0.5*(\x-1.5)*(\x-1.5) + 2.2});
\end{tikzpicture}
    \item \begin{tikzpicture}[scale=0.5, baseline=(current bounding box.center)]
    % Assi
    \draw[-{Stealth[scale=1.2]}, thick] (0,0) -- (3.5,0) node[below, pos=0.8] {posizione};
    \draw[-{Stealth[scale=1.2]}, thick] (0,0) -- (0,3.5) node[left, rotate=90, pos=0.8, anchor=south] {portata};
    
    % Funzione a campana (Gaussiana) su base costante
    % y = baseline + ampiezza * exp(-(x-centro)^2 / larghezza)
    \draw[thick, samples=100, domain=0:3] plot (\x, {0.6 + 1.4*exp(-(\x-1.5)*(\x-1.5)/0.2)});
\end{tikzpicture}
    \item \begin{tikzpicture}[scale=0.5, baseline=(current bounding box.center)]
    % Assi
    \draw[-{Stealth[scale=1.2]}, thick] (0,0) -- (3.5,0) node[below, pos=0.8] {posizione};
    \draw[-{Stealth[scale=1.2]}, thick] (0,0) -- (0,3.5) node[left, rotate=90, pos=0.8, anchor=south] {portata};
    
    % Funzione a "buca" invertita
    % y = baseline - ampiezza * exp(-(x-centro)^2 / larghezza)
    \draw[thick, samples=100, domain=0:3] plot (\x, {2.0 - 1.2*exp(-(\x-1.5)*(\x-1.5)/0.2)});
\end{tikzpicture}
}

\itemdomanda{Un sub risente di una pressione di 2 atm. Si può affermare che:}{
    \item è a circa 10 m di profondità
    \item è a circa 20 m di profondità
    \item è a circa 30 m di profondità
    \item è a circa 40 m di profondità
    \item sta risalendo velocemente
}

\itemdomanda{Tubo sezione costante ad U aperto un ramo e chiuso tappo T altro. Forza acqua sul tappo? \\ \begin{tikzpicture}[>=Stealth, thick, scale=0.6, transform shape, baseline=(current bounding box.center)]

    % --- DEFINIZIONE PARMETRI ---
    \def\tubeW{1.0}    % Larghezza del tubo
    \def\innerR{1.0}   % Raggio della curva interna
    % Il raggio esterno è calcolato automaticamente
    \pgfmathsetmacro{\outerR}{\innerR + \tubeW}
    
    \def\liqH{5.0}     % Altezza liquido a sinistra (H) dai centri dei cerchi
    \def\liqh{3.0}     % Altezza liquido a destra (h) dai centri dei cerchi
    \def\blockH{1.2}   % Altezza del blocco T
    
    % Coordinate di base: Il centro dei semicerchi inferiori è (0,0)

    % --- 1. LIQUIDO (Riempimento) ---
    % Disegniamo il contorno del liquido e lo riempiamo
    \fill[gray!40] 
        (-\outerR, \liqH) -- (-\outerR, 0)          % Giù a sinistra esterno
        arc[start angle=180, end angle=360, radius=\outerR] % Curva esterna
        -- (\outerR, \liqh)                         % Su a destra esterno
        -- (\innerR, \liqh)                         % Attraverso superficie dx
        -- (\innerR, 0)                             % Giù a destra interno
        arc[start angle=0, end angle=-180, radius=\innerR]  % Curva interna
        -- (-\innerR, \liqH)                        % Su a sinistra interno
        -- cycle;                                   % Chiude superficie sx

    % --- 2. TUBO (Contorni) ---
    % Contorno esterno
    \draw[thick] (-\outerR, \liqH + 0.8) -- (-\outerR, 0) 
        arc[start angle=180, end angle=360, radius=\outerR] 
        -- (\outerR, \liqh + \blockH + 0.5);
        
    % Contorno interno
    \draw[thick] (-\innerR, \liqH + 0.8) -- (-\innerR, 0) 
        arc[start angle=180, end angle=360, radius=\innerR] % Corretto: ora segue il fondo
        -- (\innerR, \liqh + \blockH + 0.8);

    % --- 3. OGGETTO T ---
    \fill[black] (\innerR, \liqh) rectangle (\outerR, \liqh + \blockH);
    \node[white, font=\bfseries\large] at ($(\innerR, \liqh) + 0.5*(\tubeW, \blockH)$) {T};

    % --- 4. LIVELLI E QUOTE ---
    % Linea di riferimento inferiore (il punto più basso del tubo)
    \coordinate (BottomPoint) at (0, -\outerR);
    \draw[thin] (-\outerR - 1.2, -\outerR) -- (\outerR + 1.2, -\outerR);

    % Livello 1 (Sinistra)
    \draw[thin] (-\outerR, \liqH) -- (\outerR + 2.0, \liqH) node[right, font=\small] {livello 1};
    % Quota H
    \draw[<->, thin] (-\outerR - 0.8, -\outerR) -- (-\outerR - 0.8, \liqH) node[midway, fill=white, font=\small] {$H$};

    % Livello 2 (Destra)
    \draw[thin] (\outerR, \liqh) -- (-\outerR - 2.0, \liqh) node[left, font=\small] {livello 2};
    % Quota h
    \draw[<->, thin] (\outerR + 0.8, -\outerR) -- (\outerR + 0.8, \liqh) node[midway, fill=white, font=\small] {$h$};

    % --- 5. ELEMENTO S ---
    \node[circle, draw, thin, inner sep=1.5pt, font=\itshape\small, fill=white] at (0, -\outerR + \tubeW/2) {S};

\end{tikzpicture} \\ }{
    \item $\rho g (H-h)$ verso il basso
    \item $\rho g h$ verso l'alto
    \item $\rho g H S$ verso l'alto
    \item $\rho g (H+h)$ verso il basso
    \item $\rho g (H-h) S$ verso l'alto
}

\itemdomanda{Un oggetto è posto a una distanza di 60 cm da una lente sottile convergente. L'immagine prodotta dalla lente risulta capovolta e di dimensioni pari alla metà di quelle dell'oggetto originale. Qual è la distanza focale della lente?}{
    \item 60 cm
    \item 45 cm
    \item 30 cm
    \item 20 cm
    \item 90 cm
}

\itemdomanda{Strisce colorate macchie olio pozzanghere dovute a:}{
    \item interferenza e diffrazione
    \item riflettività acqua/olio
    \item cielo diffonde colori
    \item interferenza strato sottile olio
    \item diffrazione luce
}

\itemdomanda{Perché forchetta legno è meglio della metallica per rigirare spaghetti?}{
    \item conducibilità termica metalli molto più grande
    \item calore specifico metalli maggiore
    \item conducibilità elettrica metalli maggiore
    \item legno si elettrizza meno
    \item legno è più leggero
}

\itemdomanda{Stesso calore a vetro e rame, $c_{vetro} > c_{rame}$. Allora:}{
    \item rame riscalda prima ma T finale minore
    \item rame avrà temperatura finale maggiore
    \item vetro avrà temperatura finale maggiore
    \item rame riscalda prima ma T finale uguale
    \item temperature finali uguali
}

\newpage
\section*{Appendice: risposte corrette}
\begin{longtable}{@{}r>{\raggedright\arraybackslash}p{0.14\linewidth}>{\raggedright\arraybackslash}p{0.20\linewidth}>{\raggedright\arraybackslash}p{0.42\linewidth}c@{}}
\toprule
\textbf{N.} & \textbf{Materia} & \textbf{Argomento} & \textbf{Titolo} & \textbf{Risp.} \\
\midrule
\endfirsthead
\toprule
\textbf{N.} & \textbf{Materia} & \textbf{Argomento} & \textbf{Titolo} & \textbf{Risp.} \\
\midrule
\endhead
\midrule
\multicolumn{5}{r}{\footnotesize\textit{Continua}} \\
\endfoot
\bottomrule
\endlastfoot
1 & Logica & Logica Analitica & Giorni della settimana & A \\
2 & Logica & Logica Analitica & Ordinamento pesi & B \\
3 & Logica & Logica Analitica & Posizionamento lineare & D \\
4 & Logica & Logica Analitica & Relazioni d'ordine & E \\
5 & Logica & Logica Numerica & Inserimento in griglia & A \\
6 & Logica & Logica Numerica & Successione quadratica & C \\
7 & Scienze & Fluidi & Portata massa grafico & B \\
8 & Scienze & Fluidi & Profondità subacquea & A \\
9 & Scienze & Fluidi & Tubo sezione U & E \\
10 & Scienze & Onde e Ottica & Focale lente convergente & D \\
11 & Scienze & Onde e Ottica & Macchie olio pozzanghere & D \\
12 & Scienze & Termodinamica & Forchetta legno spaghetti & A \\
13 & Scienze & Termodinamica & Vetro e rame & B
\end{longtable}

\end{document}
